\documentclass[conference,onecolumn]{IEEEtran}
\IEEEoverridecommandlockouts
\usepackage{cite}
\usepackage{amsmath,amssymb,amsfonts}
\usepackage{algorithmic}
\usepackage{graphicx}
\usepackage{textcomp}
\usepackage{xcolor}
\usepackage{hyperref}
\newcommand{\srclink}[2]{\href{#1}{\textcolor{blue}{#2}}}
\usepackage[ruled,vlined]{algorithm2e}
\usepackage{booktabs}
\usepackage{multirow}
\usepackage[hmargin=2.5cm]{geometry} % side margins (text width = page width - 5 cm)

\def\BibTeX{{\rm B\kern-.05em{\sc i\kern-.025em b}\kern-.08em
    T\kern-.1667em\lower.7ex\hbox{E}\kern-.125emX}}

\setlength{\parindent}{0pt}
\usepackage{caption}
\setlength{\textfloatsep}{6pt plus 1.0pt minus 1.5pt}
\setlength{\floatsep}{6pt plus 1.0pt minus 1.5pt}
\setlength{\intextsep}{6pt plus 1.0pt minus 1.5pt}

\begin{document}

\title{Parallel Expectation Maximization Clustering\\
{\normalsize\normalfont\itshape Department of Information Engineering and Computer Science, University of Trento}\\
{\normalsize\normalfont\url{https://github.com/HPC-2627/ParallelExpectationMaximizationClustering}}}

\author{
\IEEEauthorblockN{Michael Bernasconi}
\IEEEauthorblockA{\textit{Student ID: 267681}\\
\textit{michael.bernasconi@studenti.unitn.it}}
\and
\IEEEauthorblockN{Gabriele Fronzoni}
\IEEEauthorblockA{\textit{Student ID: 267336}\\
\textit{gabriele.fronzoni@studenti.unitn.it}}
}

\maketitle

\begin{abstract}
\end{abstract}

\begin{IEEEkeywords}
\end{IEEEkeywords}

\section{Introduction}
Clustering groups the points of a dataset so that points in the same group are similar. Here we assume that the data come from a mixture of Gaussian distributions and that the group of each point is unknown; the parameters of the mixture are estimated with the Expectation--Maximization (EM) algorithm. EM is iterative and every iteration scans the whole dataset: its cost grows with the data, which makes it a good candidate for parallelization. The goal of the project is to implement EM in C with a hybrid MPI+OpenMP design, to evaluate optimizations of computation and communication, and to measure their scalability on multiple nodes (strong and weak scaling) on real and synthetic datasets.

\subsection{The EM algorithm for Gaussian mixtures}
We distinguish two things: the \emph{model} is the assumption about the shape of the data (a mixture of Gaussians, with still unknown parameters), while \emph{EM} is the iterative procedure that estimates those parameters from the data.

\paragraph{Model.} The input is a matrix $X \in \mathbb{R}^{N\times D}$ of $N$ points (rows) with $D$ features (columns). We assume each point is generated by one of $K$ Gaussian components: a hidden label $z_i \sim \mathrm{Cat}(\pi)$ selects the component, then $x_i \mid z_i{=}k \sim \mathcal{N}(\mu_k,\Sigma_k)$. Marginalizing $z_i$, the density of a point is
\begin{equation}
  p(x) = \sum_{k=1}^{K} \pi_k \, \mathcal{N}(x \mid \mu_k, \Sigma_k),
  \qquad \pi_k \ge 0,\ \sum_{k=1}^{K} \pi_k = 1,
\end{equation}
where $\pi_k$ is the mixing weight of component $k$, and $\mu_k$, $\Sigma_k$ are its mean and covariance matrix. Let $\theta=\{\pi_k,\mu_k,\Sigma_k\}_{k=1}^{K}$ be the set of all parameters. Since the labels $z_i$ are not observed, we estimate $\theta$ by maximum likelihood; for i.i.d.\ points the log-likelihood is
\begin{equation}
  \mathcal{L}(\theta) = \log p(X \mid \theta) = \sum_{i=1}^{N} \log \sum_{k=1}^{K} \pi_k \, \mathcal{N}(x_i \mid \mu_k,\Sigma_k).
\end{equation}
The logarithm of a sum prevents a closed-form maximizer: the parameters of the components cannot be separated (they would be, if the $z_i$ were known). We therefore use the EM algorithm, which iteratively increases $\mathcal{L}(\theta)$.

In the main line of this work the covariance is \emph{diagonal}, $\Sigma_k=\mathrm{diag}(\sigma_{k1}^2,\dots,\sigma_{kD}^2)$, so the density factorizes over the features:
\begin{equation}
  \mathcal{N}(x \mid \mu_k,\Sigma_k) = \prod_{d=1}^{D} \frac{1}{\sqrt{2\pi\sigma_{kd}^2}} \exp\!\left(-\frac{(x_d-\mu_{kd})^2}{2\sigma_{kd}^2}\right).
\end{equation}
This costs $O(D)$ per point and component instead of the $O(D^3)$ matrix inversion (plus $O(D^2)$ per evaluation) of a full covariance, and needs only $2D$ instead of $D+D(D+1)/2$ parameters per component.

\paragraph{E-step.} For every point $i$ and cluster $k$ it computes the \emph{responsibility} $r_{ik}\in[0,1]$, the probability that $x_i$ belongs to cluster $k$ under the current parameters ($\sum_k r_{ik}=1$); $j$ runs over all the clusters:
\begin{equation}
r_{ik}=\frac{\pi_k\,\mathcal{N}(x_i\mid\mu_k,\sigma_k^2)}{\sum_{j=1}^{K}\pi_j\,\mathcal{N}(x_i\mid\mu_j,\sigma_j^2)}.
\end{equation}

\paragraph{M-step.} It recomputes the parameters weighting each point by $r_{ik}$. With $N_k=\sum_{i=1}^{N}r_{ik}$ (effective number of points of cluster $k$):
\begin{gather}
\mu_k=\tfrac{1}{N_k}\sum_{i=1}^{N} r_{ik}\,x_i,\qquad
\pi_k=\tfrac{N_k}{N},\\
\sigma^2_{kd}=\tfrac{1}{N_k}\sum_{i=1}^{N} r_{ik}\,x_{id}^2-\mu_{kd}^2.
\end{gather}
In the code the M-step accumulates in a single pass the sums $N_k$, $\sum_i r_{ik}x_i$ and $\sum_i r_{ik}x_i^2$ (the latter component by component) and then converts them into the parameters.

\paragraph{Initialization and stopping.} The means $\mu_k$ are $K$ distinct rows of $X$ chosen with a seeded generator (so every process obtains the same rows without communicating), $\pi_k=1/K$ and $\sigma^2_{kd}=1$ (the data are standardized: mean 0 and variance 1 per variable). After every E-step $\mathcal{L}$ is computed and EM stops if $|\mathcal{L}-\mathcal{L}_{\mathrm{prev}}|\le\varepsilon|\mathcal{L}|$, where $\mathcal{L}_{\mathrm{prev}}$ is the value of the previous iteration ($-\infty$ at the start) and $\varepsilon$ a relative tolerance, or after $T_{\max}$ iterations (defaults $\varepsilon=10^{-6}$, $T_{\max}=200$). $\mathcal{L}$ never decreases~\cite[Sec.~9.4]{bishop} but only a local maximum is reached: with $R$ restarts the run with the highest $\mathcal{L}$ is kept ($R=1$ in the measurements).

\begin{algorithm}[t]
\caption{EM for a Gaussian mixture (one run)}\label{alg:em}
\small
\KwIn{$X$ ($N\times D$), $K$, $\varepsilon$, $T_{\max}$}
\KwOut{$\theta$, labels $\arg\max_k r_{ik}$}
Initialize $\theta$: $\mu_k\leftarrow$ random distinct rows, $\pi_k\leftarrow 1/K$, $\sigma^2_{kd}\leftarrow 1$\;
$\mathcal{L}_{\mathrm{prev}}\leftarrow-\infty$\;
\For{$t=1$ \KwTo $T_{\max}$}{
  E-step: compute $r_{ik}$ for all $i,k$ and $\mathcal{L}$\;
  \lIf{$|\mathcal{L}-\mathcal{L}_{\mathrm{prev}}|\le\varepsilon|\mathcal{L}|$}{break}
  M-step: compute $N_k$, $\mu_k$, $\pi_k$, $\sigma^2_{kd}$ for all $k,d$\;
  $\mathcal{L}_{\mathrm{prev}}\leftarrow\mathcal{L}$\;
}
\lIf{$T_{\max}$ reached without convergence}{repeat the E-step (so $\mathcal{L}$ and $r_{ik}$ match $\theta$)}
\end{algorithm}

\paragraph{Numerical safeguards.} The formulas, implemented literally, produce errors in finite-precision arithmetic; the code adopts three protections.
\begin{itemize}
\item Computation in log scale. With large $D$ the densities $\mathcal{N}$ are tiny numbers that the computer rounds to 0, and the normalization of $r_{ik}$ would give $0/0$. For each point we therefore compute the logarithms $v_k=\ln\big(\pi_k\,\mathcal{N}(x_i\mid\mu_k,\sigma_k^2)\big)$ and sum them with
$\ln\sum_{k}e^{v_k}=m+\ln\sum_{k}e^{v_k-m}$, where $m=\max_k v_k$: the largest term becomes $e^0=1$ and the sum is never zero. The result is $\ln p(x_i)$, which contributes to $\mathcal{L}$ and gives $r_{ik}=e^{v_k-\ln p(x_i)}$.
\item Non-zero variances. If a cluster shrinks to a few almost identical points, $\sigma^2_{kd}$ tends to 0 and the density to infinity, making $\mathcal{L}$ unstable~\cite[Sec.~9.2.1]{bishop}. We therefore enforce $\sigma^2_{kd}\ge\delta=10^{-6}$ at every M-step.
\item Empty clusters. If $N_k<10^{-10}$ no point is assigned to the cluster and the formulas would divide by zero: mean and variances stay as before and the weight is set to the minimum $\pi_k=10^{-12}$, never 0, because the E-step uses $\ln\pi_k$.
\end{itemize}

\subsection{Sequential implementations and cost}
We wrote two \emph{sequential} versions, \emph{only to understand the algorithm in depth}: they are not the performance baseline, which is the diagonal MPI+OpenMP version, against which the diagonal optimizations (described in the following sections) are evaluated; optimization~1 of the full model is evaluated against the full MPI+OpenMP baseline.
\begin{itemize}
\item Diagonal: the model described above, with one variance $\sigma^2_{kd}$ per variable.
\item Full: each cluster has a complete covariance matrix $\Sigma_k\in\mathbb{R}^{D\times D}$, which also models the correlations between variables. The model changes, so its results are compared only with other full-covariance runs. The algorithm does not change (Algorithm~\ref{alg:em} is the same, with the same safeguards): only $\sigma^2_{kd}$ is replaced by $\Sigma_k$ in the density of the E-step and in the covariance update of the M-step. Besides the sequential version, the full model also exists in MPI+OpenMP form: baseline and optimization~1.
\end{itemize}

\paragraph{Cost of one iteration.} With $N$ points, $K$ clusters and $D$ variables:
\begin{itemize}
\item Diagonal, $O(NKD)$. For each of the $NK$ pairs $(i,k)$ the E-step computes $m_k(x_i)$ summing $D$ terms, and the M-step accumulates $D$ terms in $\sum_i r_{ik}x_i$ and $D$ in $\sum_i r_{ik}x_i^2$.
\item Full, $O(NKD^2)$. For each pair $(i,k)$ the Mahalanobis distance and the accumulation of $\Sigma_k$ cost $O(D^2)$ instead of $O(D)$. Factorizing the $K$ matrices $\Sigma_k$ adds $O(KD^3)$, which does not depend on $N$ and is negligible if $N\gg D$.
\item Common, $O(NK)$. The normalization of the responsibilities and the computation of $\mathcal{L}$ cost $K$ operations per point.
\end{itemize}

\section{Parallel Design Methodology}\label{sec:design}
\subsection{Foster's Design Strategy}
\subsection{Communication}

\section{Experimental Setup \& Dataset}
\subsection{Hardware Infrastructure}
\subsection{Dataset}

\begin{table*}[p]
\centering
\caption{Real datasets. $N$ and $D$ are the values actually used by the EM after preprocessing; $K$ is the number of ground-truth classes (used only for evaluation).}
\label{tab:real}
\footnotesize
\setlength{\tabcolsep}{3pt}
\resizebox{\textwidth}{!}{%
\begin{tabular}{llrrrrll}
\toprule
\textbf{Name} & \textbf{Source} & \textbf{$N$} & \textbf{$D$} & \textbf{$K$} & \textbf{File} & \textbf{Label} & \textbf{Description} \\
\midrule
\texttt{statlog\_heart}   & \srclink{https://archive.ics.uci.edu/dataset/145}{UCI \#145} & 270 & 13 & 2 & 28 KB & \texttt{heart-disease} & Heart disease, small clinical sample \\
\texttt{heart\_failure}   & \srclink{https://archive.ics.uci.edu/dataset/519}{UCI \#519} & 299 & 12 & 2 & 29 KB & \texttt{death\_event} & Heart failure patients, survival \\
\texttt{gene\_rnaseq}     & \srclink{https://archive.ics.uci.edu/dataset/401}{UCI \#401} & 801 & 20\,264 & 5 & 130 MB & \texttt{Class} & TCGA RNA-Seq, 5 cancer types \\
\texttt{cardiotocography} & \srclink{https://archive.ics.uci.edu/dataset/193}{UCI \#193} & 2\,126 & 21 & 10 & 357 KB & \texttt{CLASS} & Fetal cardiotocograms \\
\texttt{parkinson}        & \srclink{https://archive.ics.uci.edu/dataset/189}{UCI \#189} & 5\,875 & 19 & -- & 893 KB & -- (regression) & Voice telemonitoring \\
\texttt{eeg\_eye\_state}  & \srclink{https://archive.ics.uci.edu/dataset/264}{UCI \#264} & 14\,980 & 14 & 2 & 1.7 MB & \texttt{eyeDetection} & EEG, eyes open/closed \\
\texttt{cardiovascular}   & \srclink{https://www.openml.org/d/45547}{OpenML 45547} & 70\,000 & 11 & 2 & 6.2 MB & \texttt{cardio} & Cardiovascular disease \\
\texttt{diabetes130}      & \srclink{https://archive.ics.uci.edu/dataset/296}{UCI \#296} & 101\,766 & 11 & 3 & 9.0 MB & \texttt{readmitted} & Hospital readmission, 130 US hospitals \\
\texttt{cdc\_diabetes}    & \srclink{https://archive.ics.uci.edu/dataset/891}{UCI \#891} & 253\,680 & 21 & 2 & 43 MB & \texttt{Diabetes\_binary} & CDC health indicators survey \\
\texttt{brainweb}         & \srclink{https://brainweb.bic.mni.mcgill.ca/}{BrainWeb (subj.~44)} & 11\,862\,016 & 1 & 12 & 95 MB & crisp labels & Simulated T1 brain MRI, one voxel = one point \\
\bottomrule
\end{tabular}}
\end{table*}
\clearpage

\begin{table*}[p]
\centering
\caption{Synthetic profiles. $\delta$ is the distance between cluster means in units of the within-cluster standard deviation.}
\label{tab:synth-profiles}
\footnotesize
\setlength{\tabcolsep}{3pt}
\begin{tabular}{llll}
\toprule
 & \texttt{synth\_clean} & \texttt{synth\_overlap} & \texttt{synth\_rare} \\
\midrule
Role & easy, satisfies the model & realistic & hard \\
Clusters $K$ & 4 & 5 & 6 \\
Variables $D$ & 10 & 20 & 50 \\
Cluster weights & 25, 25, 25, 25\,\% & 30, 25, 20, 15, 10\,\% & 40, 25, 18, 12, 4, 1\,\% \\
Separation $\delta$ & 6.0 (well separated) & 2.5 (overlapping) & 4.0 (intermediate) \\
Covariance in a cluster & diagonal & correlated (3-factor model, 0.5) & diagonal \\
Std.\ per variable & 1.0 & 1.0 (before correlation) & in $[0.6, 1.4]$ \\
Outliers & 0\,\% & 0\,\% & 2\,\% (label $-1$) \\
Skewed (log-normal) variables & no & yes (7 of 20) & no \\
Physiological clipping & no & yes & yes (clinical variables) \\
What it tests & correctness, scaling reference & slow convergence, choice of $K$, model mismatch & numerical stability \\
\bottomrule
\end{tabular}
\end{table*}
\clearpage

\begin{table*}[p]
\centering
\caption{Clinical variables used to build the synthetic data (typical mean $\pm$ std and physiological limits).}
\label{tab:synth-vars}
\footnotesize
\begin{tabular}{rlllcc}
\toprule
\# & \textbf{Variable} & \textbf{Unit} & \textbf{Mean $\pm$ std} & \textbf{Limits} & \textbf{Skew} \\
\midrule
1 & \texttt{age} & years & $55 \pm 15$ & 18--100 & \\
2 & \texttt{bmi} & kg/m$^2$ & $27 \pm 5$ & 14--60 & \\
3 & \texttt{sys\_bp} & mmHg & $130 \pm 20$ & 80--230 & \\
4 & \texttt{dia\_bp} & mmHg & $80 \pm 11$ & 40--130 & \\
5 & \texttt{glucose} & mg/dL & $105 \pm 30$ & 50--400 & \checkmark \\
6 & \texttt{hba1c} & \% & $5.8 \pm 1.0$ & 4--14 & \checkmark \\
7 & \texttt{cholesterol} & mg/dL & $200 \pm 40$ & 90--400 & \\
8 & \texttt{hdl} & mg/dL & $52 \pm 14$ & 20--110 & \\
9 & \texttt{creatinine} & mg/dL & $1.0 \pm 0.3$ & 0.3--6 & \checkmark \\
10 & \texttt{heart\_rate} & bpm & $75 \pm 12$ & 40--150 & \\
11 & \texttt{ldl} & mg/dL & $120 \pm 35$ & 30--300 & \\
12 & \texttt{triglycerides} & mg/dL & $140 \pm 70$ & 30--800 & \checkmark \\
13 & \texttt{albumin} & g/dL & $4.2 \pm 0.4$ & 2--5.5 & \\
14 & \texttt{hemoglobin} & g/dL & $13.5 \pm 1.8$ & 6--19 & \\
15 & \texttt{wbc} & $10^3/\mu$L & $7.0 \pm 2.0$ & 2--25 & \checkmark \\
16 & \texttt{platelets} & $10^3/\mu$L & $250 \pm 60$ & 50--600 & \\
17 & \texttt{sodium} & mmol/L & $140 \pm 3$ & 120--160 & \\
18 & \texttt{potassium} & mmol/L & $4.2 \pm 0.4$ & 2.5--6.5 & \\
19 & \texttt{alt} & U/L & $28 \pm 15$ & 5--200 & \checkmark \\
20 & \texttt{crp} & mg/L & $3.0 \pm 3.0$ & 0.1--100 & \checkmark \\
21--50 & \texttt{biomarker\_01..30} & -- & $0 \pm 1$ & none & \\
\bottomrule
\end{tabular}
\end{table*}
\clearpage

\begin{table*}[p]
\centering
\caption{Weak scaling: rows and data file size for each profile ($N_0 = 250\,000$ rows per process; labels file of $4N$ bytes for every profile). MB $= 10^6$ bytes.}
\label{tab:weak}
\footnotesize
\setlength{\tabcolsep}{3pt}
\begin{tabular}{rrrrrr}
\toprule
\textbf{$p$} & \textbf{$N = N_0 \cdot p$} & \texttt{synth\_clean} & \texttt{synth\_overlap} & \texttt{synth\_rare} & \textbf{Labels} \\
\midrule
1  & 250\,000   & 20 MB  & 40 MB  & 100 MB   & 1 MB \\
2  & 500\,000   & 40 MB  & 80 MB  & 200 MB   & 2 MB \\
4  & 1\,000\,000 & 80 MB  & 160 MB & 400 MB   & 4 MB \\
8  & 2\,000\,000 & 160 MB & 320 MB & 800 MB   & 8 MB \\
16 & 4\,000\,000 & 320 MB & 640 MB & 1\,600 MB & 16 MB \\
\midrule
\textbf{Total} & & 620 MB & 1\,240 MB & 3\,100 MB & 31 MB \\
\bottomrule
\end{tabular}
\end{table*}
\clearpage

\begin{table*}[p]
\centering
\caption{Strong scaling (fixed $N = 4\,000\,000$): data per rank and work per rank $N \cdot K \cdot D / p$. The work is an operation count proportional to the cost of an iteration, not a measured time.}
\label{tab:strong}
\footnotesize
\setlength{\tabcolsep}{3pt}
\begin{tabular}{rrrrrrrr}
\toprule
 & & \multicolumn{3}{c}{\textbf{Data per rank}} & \multicolumn{3}{c}{\textbf{Work per rank ($N K D / p$)}} \\
\cmidrule(lr){3-5} \cmidrule(lr){6-8}
\textbf{$p$} & \textbf{Rows/rank} & \texttt{clean} & \texttt{overlap} & \texttt{rare} & \texttt{clean} & \texttt{overlap} & \texttt{rare} \\
\midrule
1  & 4\,000\,000 & 320 MB & 640 MB & 1\,600 MB & $1.6 \cdot 10^8$ & $4.0 \cdot 10^8$ & $1.2 \cdot 10^9$ \\
2  & 2\,000\,000 & 160 MB & 320 MB & 800 MB    & $8.0 \cdot 10^7$ & $2.0 \cdot 10^8$ & $6.0 \cdot 10^8$ \\
4  & 1\,000\,000 & 80 MB  & 160 MB & 400 MB    & $4.0 \cdot 10^7$ & $1.0 \cdot 10^8$ & $3.0 \cdot 10^8$ \\
8  & 500\,000    & 40 MB  & 80 MB  & 200 MB    & $2.0 \cdot 10^7$ & $5.0 \cdot 10^7$ & $1.5 \cdot 10^8$ \\
16 & 250\,000    & 20 MB  & 40 MB  & 100 MB    & $1.0 \cdot 10^7$ & $2.5 \cdot 10^7$ & $7.5 \cdot 10^7$ \\
\bottomrule
\end{tabular}
\end{table*}
\clearpage

\section{Performance Evaluation}
\subsection{Execution Time and Throughput}
\subsection{Strong Scaling}
\subsection{Weak Scaling}
\subsection{Communication vs Computation Breakdown}

\section{Discussion}

\section{Conclusion}

\clearpage
\begin{thebibliography}{1}
\bibitem{bishop} C. M. Bishop, \href{https://www.microsoft.com/en-us/research/uploads/prod/2006/01/Bishop-Pattern-Recognition-and-Machine-Learning-2006.pdf}{\textcolor{blue}{\emph{Pattern Recognition and Machine Learning}}}, Springer, 2006, ch.~9.
\end{thebibliography}

\end{document}
